% sec_experiments.tex — Section 5: Experiments and Simulation Results
% ICGNC 2026, Springer svproc template

\section{Experiments and Simulation Results}
\label{sec:experiments}

\subsection{Simulation Setup}

We simulate a partially observable grid battlefield with local $11\!\times\!11$
patches (obs radius $=5$) and shared global rewards as in Sect.~\ref{sec:problem}.
The default full-scale configuration follows
\texttt{configs/default.yaml}: a $100\!\times\!100$ grid, $10$~UAVs,
$5$~targets, and $7$~threat zones, with communication radius
$R_{\mathrm{comm}}=20$ and $\texttt{max\_steps}=200$ per episode.

\textbf{TAGA-MAPPO training.} We adopt three-stage curriculum learning:
\begin{itemize}\setlength{\itemsep}{0pt}\setlength{\parskip}{0pt}
  \item \texttt{stage\_1\_easy}\,: $3$~UAVs, $2$~targets, $0$~threats, $30\!\times\!30$;
  \item \texttt{stage\_2\_medium}\,: $6$~UAVs, $3$~targets, $3$~threats, $50\!\times\!50$;
  \item \texttt{stage\_3\_hard}\,: $10$~UAVs, $5$~targets, $7$~threats, $100\!\times\!100$.
\end{itemize}
\noindent MAPPO uses GAE($\lambda=0.95$), PPO clip
$\varepsilon=0.2$, discount $\gamma=0.99$, Adam learning rate $3\times10^{-4}$,
and entropy coefficient $0.05$ for the short sanity runs reported below.

\textbf{Baselines for the quick validation run.}  We compare (i)~a \emph{random
joint policy} that uniformly samples hierarchical actions on the
\texttt{stage\_3\_hard} layout, and (ii)~a \emph{greedy TAGA-MAPPO policy}
(argmax over mode and move logits) after a \emph{limited} training budget of
$90$~PPO updates on CPU, intended as an engineering sanity check rather than a
fully converged benchmark.  Full-scale comparisons against classical planners
(e.g., ACO), flat MAPPO, and TAGA-LLM under five directives
(\textsc{bal}, \textsc{sch}, \textsc{atk}, \textsc{evd}, \textsc{coo}) are left
to the extended evaluation track with $30$~seeds and Welch's $t$-test at
$\alpha=0.05$, as originally planned.

\textbf{Hardware:} $8\!\times\!$A100 GPU server; Stage~I training ${\sim}10$\,h,
Stage~II LoRA fine-tuning ${\sim}37$\,min (single GPU).

\subsection{Simulation Results}

\begin{table}[t]
\centering
\caption{Rapid local validation (mean $\pm$ std).  Random: 24 episodes on
\texttt{stage\_3\_hard}.  TAGA-MAPPO: greedy rollouts after 90 updates
(10 episodes on the post-training environment layout).}
\label{tab:sanity}
\begin{tabular}{lcc}
\toprule
Metric & Random & TAGA-MAPPO (greedy, short train) \\
\midrule
Search coverage      & $0.361 \pm 0.043$ & $0.557 \pm 0.068$ \\
Target destroy rate  & $0.183 \pm 0.143$ & $0.000 \pm 0.000$ \\
Norm.\ completion time & $0.900 \pm 0.000$ & $1.000 \pm 0.000$ \\
UAV survival rate    & $0.742 \pm 0.125$ & $0.620 \pm 0.132$ \\
Coop.\ attack ratio  & $0.000 \pm 0.000$ & $0.000 \pm 0.000$ \\
\bottomrule
\end{tabular}
\end{table}

Table~\ref{tab:sanity} shows that even with a small optimization budget,
TAGA-MAPPO already improves search coverage relative to random exploration;
the destroy-rate and time-to-completion metrics indicate that substantially
longer training (and LLM distillation) are required before claiming mission
completion at scale---consistent with the staged experimental roadmap.

\subsection{Instruction Controllability}

We will evaluate TAGA-LLM under five operational directives
(\textsc{bal}, \textsc{sch}, \textsc{atk}, \textsc{evd}, \textsc{coo}) using
the topology-aware prompt encoder described in Sect.~\ref{sec:distillation}.
The distilled model is expected to shift coverage--aggression trade-offs
measurably across directives; quantitative curves are omitted here pending
completion of the LoRA fine-tuning pipeline.

\subsection{Ablation Studies}

Planned ablations include (i)~TAGA-aware prompt encoding versus a flat text
serialization; (ii)~chain-of-thought supervision versus action-only supervision;
(iii)~counterfactual data augmentation on/off; (iv)~ablating individual TAGA
factors ($\lambda_d, \lambda_m, \lambda_t$) in the expert policy.  Each
ablation will be reported with the same five metrics as the main table once the
long-run training budget is finalized.

\subsection{Scalability and Robustness}

We target scalability sweeps with $5/10/15/20$~UAVs, communication dropout at
$10\%/20\%/30\%$, and inference latency benchmarks for the distilled LLM.
These experiments are not included in the rapid CPU sanity run above.

% TODO: 图片文件缺失，暂时注释
% \begin{figure}[t]
% \centering
% \includegraphics[width=\textwidth]{training_curves.pdf}
% \caption{TAGA-MAPPO training traces under curriculum scheduling
% (policy/value loss, entropy, reward/coverage).  Vertical markers indicate
% automatic stage transitions when rolling coverage exceeds stage thresholds.}
% \label{fig:training_curves}
% \end{figure}
%
% \begin{figure}[t]
% \centering
% \includegraphics[width=0.6\textwidth]{curriculum_stage_timeline.pdf}
% \caption{Curriculum stage occupancy over PPO updates for the same sanity run.}
% \label{fig:curriculum_timeline}
% \end{figure}
